\documentclass[11pt]{article}

% Change "review" to "final" to generate the final (sometimes called camera-ready) version.
% Change to "preprint" to generate a non-anonymous version with page numbers.
\usepackage[preprint]{acl}

% Standard package includes
\usepackage{times}
\usepackage{latexsym}
\usepackage{amsmath}
\usepackage[most]{tcolorbox}
\usepackage{booktabs}
% For proper rendering and hyphenation of words containing Latin characters (including in bib files)
\usepackage[T1]{fontenc}
% For Vietnamese characters
% \usepackage[T5]{fontenc}
% See https://www.latex-project.org/help/documentation/encguide.pdf for other character sets

% This assumes your files are encoded as UTF8
\usepackage[utf8]{inputenc}
\usepackage{url}
% This is not strictly necessary, and may be commented out,
% but it will improve the layout of the manuscript,
% and will typically save some space.
\usepackage{microtype}

% This is also not strictly necessary, and may be commented out.
% However, it will improve the aesthetics of text in
% the typewriter font.
\usepackage{inconsolata}

%Including images in your LaTeX document requires adding
%additional package(s)
\usepackage{graphicx}

% If the title and author information does not fit in the area allocated, uncomment the following
%
%\setlength\titlebox{<dim>}
%
% and set <dim> to something 5cm or larger.

\title{Rating the Raters: Rasch Measurement Theory for LLM Evaluation}

% Author information can be set in various styles:
% For several authors from the same institution:
% \author{Author 1 \and ... \and Author n \\
%         Address line \\ ... \\ Address line}
% if the names do not fit well on one line use
%         Author 1 \\ {\bf Author 2} \\ ... \\ {\bf Author n} \\
% For authors from different institutions:
% \author{Author 1 \\ Address line \\  ... \\ Address line
%         \And  ... \And
%         Author n \\ Address line \\ ... \\ Address line}
% To start a separate ``row'' of authors use \AND, as in
% \author{Author 1 \\ Address line \\  ... \\ Address line
%         \AND
%         Author 2 \\ Address line \\ ... \\ Address line \And
%         Author 3 \\ Address line \\ ... \\ Address line}

\author{Pratik Sachdeva \\
  D-Lab, University of California, Berkeley \\
  \texttt{pratik.sachdeva@berkeley.edu} \\\And
  Nathan Boudol \\
  Grenoble INP \\}

%\author{
%  \textbf{First Author\textsuperscript{1}},
%  \textbf{Second Author\textsuperscript{1,2}},
%  \textbf{Third T. Author\textsuperscript{1}},
%  \textbf{Fourth Author\textsuperscript{1}},
%\\
%  \textbf{Fifth Author\textsuperscript{1,2}},
%  \textbf{Sixth Author\textsuperscript{1}},
%  \textbf{Seventh Author\textsuperscript{1}},
%  \textbf{Eighth Author \textsuperscript{1,2,3,4}},
%\\
%  \textbf{Ninth Author\textsuperscript{1}},
%  \textbf{Tenth Author\textsuperscript{1}},
%  \textbf{Eleventh E. Author\textsuperscript{1,2,3,4,5}},
%  \textbf{Twelfth Author\textsuperscript{1}},
%\\
%  \textbf{Thirteenth Author\textsuperscript{3}},
%  \textbf{Fourteenth F. Author\textsuperscript{2,4}},
%  \textbf{Fifteenth Author\textsuperscript{1}},
%  \textbf{Sixteenth Author\textsuperscript{1}},
%\\
%  \textbf{Seventeenth S. Author\textsuperscript{4,5}},
%  \textbf{Eighteenth Author\textsuperscript{3,4}},
%  \textbf{Nineteenth N. Author\textsuperscript{2,5}},
%  \textbf{Twentieth Author\textsuperscript{1}}
%\\
%\\
%  \textsuperscript{1}Affiliation 1,
%  \textsuperscript{2}Affiliation 2,
%  \textsuperscript{3}Affiliation 3,
%  \textsuperscript{4}Affiliation 4,
%  \textsuperscript{5}Affiliation 5
%\\
%  \small{
%    \textbf{Correspondence:} \href{mailto:email@domain}{email@domain}
%  }
%}

\begin{document}
\maketitle
\begin{abstract}
LLMs now sit on every side of evaluation: as \textit{examinees} scored on benchmarks, \textit{judges} of other models' outputs, and \textit{raters} of human-generated content. Each paradigm can be viewed as a \textit{measurement} problem, where a latent property of an object is probed with items from an instrument (e.g., benchmark) by raters. Standard evaluation practices often neglect the contributions of each core component to the end result, limiting our understanding of what is being measured. Rasch measurement theory (RMT) is well-suited to this kind of problem. RMT decomposes ordinal ratings into separable \textit{facets} on a common scale. It further provides a battery of diagnostics that can identify miscalibrated measurements and rater biases. We present a case study of RMT applied to the LLM-as-rater paradigm using the \textit{Measuring Hate Speech} corpus, whose construct was itself built under RMT. We fit a series of many-facet Rasch models to annotations from nine LLMs spanning families and capability levels. Our analyses show that LLMs systematically differ from human raters in severity, item-level calibration, question-order robustness, target-identity sensitivity, and rating scale use, which all would be obscured by standard evaluation practice. Overall, we argue that RMT belongs in the toolkit for evaluating LLM-as-examinee, -judge, and -rater paradigms.\footnote{Code available at \url{https://anonymous.4open.science/r/rating_the_raters-BC63/}}
\end{abstract}

\section{Introduction}
LLMs now sit on every side of evaluation. As \textit{examinees}, they are scored on benchmarks that estimate their capabilities and character traits \cite{hendrycks2021measuring, srivastava2023beyond}. As \textit{judges}, they evaluate the outputs of other LLMs against rubrics \cite{zheng2023judging, li2024llms}. As \textit{raters}, they annotate human-generated content at scale for subjective properties like sentiment or hatefulness \cite{gilardi2023chatgpt, calderon2025alternative}. Each of these three paradigms is fundamentally a measurement problem: a latent property -- a model's reasoning ability, a response's helpfulness, a comment's hatefulness -- is probed by an instrument, scored by a rater, and reduced to a quantity on which downstream interpretations or decisions are made \cite{Wallach2025GenAIMeasurement, Weidinger2025EvaluationScience, Zhou2026GeneralScales}.

Yet evaluation practice rarely treats these problems with the tools that measurement requires. Standard reporting (accuracies, annotator agreements, etc.) collapses the object of measurement, the item, and the rater into a single aggregate score \cite{Salaudeen2025MeasurementToMeaning}. Recent work applying item response theory (IRT) and psychometric methods to NLP \cite{lalor2016building, vania2021comparing, Federiakin2025Leaderboards, Chen2025CompactIRT, Yao2025JEIRT, Ye2025LLMPsychometrics} has begun to remedy this for the LLM-as-examinee paradigm, where items can be characterized by difficulty and models by ability. The LLM-as-judge and LLM-as-rater paradigms are more recently gaining such methodologies \cite{giorgi2025human, Choi2026JudgeIRT, Feng2026NoisyJudges, Chen2026NoisyJudgeInference, Jiao2026RaterEffects}.

\textit{Rasch measurement theory} (RMT) is well-suited for this gap \cite{rasch1993probabilistic, andrich1988rasch, linacre1989manyfacet}. RMT decomposes ordinal ratings (e.g., Likert-style rubrics) into separable contributions from the object being measured, the items, and the raters, all placed on a common measurement scale. It additionally commits to two requirements that have no counterpart in standard ML evaluation: \textit{invariance}, which demands that the instrument function the same way across raters and items, and the principle that the data must fit the model rather than the model fit the data. RMT thus comes equipped with a battery of diagnostics -- rater severity, item difficulty, differential rater functioning, and rating scale functioning -- that turn rating data into structured information about each component of the measurement.

In this work, we present a case study of RMT applied to the LLM-as-rater paradigm, focusing on hate speech measurement. We use this domain because (i) hate speech is highly subjective and relevant for the LLM-as-rater paradigm, and (ii) the \textit{Measuring Hate Speech} (MHS) corpus -- 50{,}070 comments labeled by 11{,}143 annotators on a ten-item instrument -- was itself constructed under RMT, giving us a validated construct against which to compare LLM annotations \cite{DBLP:journals/corr/abs-2009-10277, sachdeva-etal-2022-measuring}. We collect annotations from nine LLMs spanning families, providers, and capability levels, and fit a series of many-facet Rasch measurement models that produce a battery of diagnostics. Across our analyses, LLMs depart from human raters in structured and interpretable ways, which would not be detectable without RMT.

Our goal is methodological: we offer a case study of how RMT supplies a principled vocabulary for interrogating LLM-as-raters, and we argue that the same vocabulary belongs in the toolkit used to evaluate LLMs across the examinee, judge, and rater paradigms. In this paper, Section~\ref{sec:rmt} provides an overview of RMT and the many-facet model. Section~\ref{sec:methods} describes the data, LLMs, and models we use. Section~\ref{sec:results} presents results in two parts: measurement scales anchored to a human baseline (Sections 4.1-4.3), and an LLM-only measurement scale that asks what the construct looks like when LLMs set the scale themselves (Section 4.4).

\section{Rasch Measurement Theory}
\label{sec:rmt}
The goal of measurement theory is to make observations of some property meaningfully comparable. In measurement settings, observations are typically produced by the interaction of three groups: the objects whose latent property is being measured, the items of a survey instrument that probe that property, and the raters who use the instrument to assign ordinal responses. Specific measurement frameworks allow one to transform these observations into variables that reflect an underlying scale, facilitating cross-rater comparisons. Item response theory (IRT) is one approach that utilizes probabilistic models to extract continuous scales from ordinal responses such as Likert ratings. 

Rasch measurement theory (RMT) is a distinct methodological tradition that commits to a specific requirement: invariance \cite{rasch1993probabilistic, andrich1988rasch}. Invariance requires that the measurement instrument function the same way regardless of which raters produce the responses, which objects are rated, and which items are used to construct the scale. RMT achieves invariance in two ways. First, it uses a specific class of additive IRT models -- here, we focus on the \textit{many-facet model} -- that mathematically guarantee invariance \cite{linacre1989manyfacet}. Second, it requires that the observational data fit the model rather than the model fit the data \cite{bond2015applying}. Thus, a poor fit should be read as information about the items, the raters, or the construct, but not as a problem to be solved by adjusting the model.

\textbf{Construct Theorization.} Constructing a Rasch measurement scale begins with theorizing a \textit{construct}: an explicit articulation of the latent property being measured and how it varies along a continuum \cite{Wallach2025GenAIMeasurement}. For hate speech, the construct specifies what it means for a comment to lie at successive points along the spectrum: from neutral or supportive speech, through expressions of disrespect, insult, humiliation, and dehumanization, to incitement of violence and calls for genocide. Developing a construct requires careful qualitative analysis of constituent examples; it is operationalized via a \textit{survey instrument} whose items map onto distinct regions of the spectrum. We adopt the construct established by prior work on the Measuring Hate Speech corpus, whose survey instrument contains ten items discussed further in Section~\ref{sec:corpus} \cite{DBLP:journals/corr/abs-2009-10277, sachdeva-etal-2022-measuring}.

\textbf{Many-Facet Rasch Model.} The many-facet model recovers a continuous scale from the ordinal responses on the ten items of the MHS instrument. It captures the decision to opt for response $k$ (say, ``strongly agree'') versus response $k-1$ (``agree''). Let $p_{nijk}$ be the probability that rater $j$ assigns comment $n$ a response $k$ on survey item $i$; similarly define $p_{nij(k-1)}$ for rating $k-1$. The model defines an \textit{odds ratio} $r_{nijk}$ as:
\begin{equation}
    \log r_{nijk} = \log \frac{p_{nijk}}{p_{nij(k-1)}} = \theta_n - \delta_i - \alpha_j - \tau_{k}. \label{eqn:base_irt}
\end{equation}
Items are coded so that higher $k$ always corresponds to greater hatefulness. A larger log-odds, therefore, indicates that the rater is more inclined to assign the more-hateful category. It depends on the following four additively separable \textit{facets}:
\begin{itemize}
    \item $\theta_n$, the \textbf{hate speech score} of comment $n$. Higher $\theta_n$ indicates a more inherently hateful comment. 
    \item $\delta_i$, the \textbf{difficulty} of survey item $i$. Difficulties provide a check on construct validity: we should expect the fitted values to follow the ordering of the construct.
    \item $\alpha_j$, the \textbf{severity} of rater $j$. Raters with higher severity are less likely to label comments as possessing features of hate speech: their thresholds for ``hatefulness'' are higher.
    \item $\tau_k$, the \textbf{Rasch-Andrich threshold} for response $k$, indicating whether the ordinal responses function appropriately.
\end{itemize}
Because comments, items, and raters share a common logit scale, model parameters can be interpreted together in the context of the construct.

\begin{figure*}[t]
    \centering
    \includegraphics[width=0.8\linewidth]{figures/figure1_model_severities.pdf}
    \caption {\textbf{Measurement models quantify LLM judge thresholds.} \textbf{a.} Fitted severity parameters $\alpha_j^{\text{LLM}}$ for nine LLMs (horizontal bars; error bars indicate standard error). Top plot compares parameter estimates (dashed colored lines) to the distribution of human severity parameters (blue curve: kernel density estimate). \textbf{b.} Comparison of severity estimates ($y$-axis) to AA Intelligence Index ($x$-axis) across expanded model set (Pearson $r=0.52$ $p<10^{-3}$). Black line denotes linear fit.  Bar and point colors denote model provider (legend shown panel \textbf{b}). }
    \label{fig:model_severities}
\end{figure*}


\section{Methods}
\label{sec:methods}
\subsection{The \textit{Measuring Hate Speech} Corpus}
\label{sec:corpus}
The \textit{Measuring Hate Speech} corpus consists of 50,070 social media comments labeled by 11,143 annotators using a survey instrument with ten items spanning the hate speech spectrum: \textit{Sentiment}, \textit{Respect}, \textit{Insult}, \textit{Humiliate}, \textit{Status}, \textit{Attack/Defend}, \textit{Dehumanize}, \textit{Violence}, \textit{Genocide}, and an overall \textit{Hate Speech} item. Each item consisted of 5 Likert-like responses, other than \textit{Hate Speech}, which was ternary. The corpus additionally contains annotations denoting the identities targeted by each comment. We focus on identity categories spanning gender, race/ethnicity, religion, and sexuality. 

We obtained the MHS corpus from HuggingFace in compliance with the authors' license and intended use \cite{ucberkeleyDlabMHS}. We consider two subsets of the MHS dataset for obtaining LLM ratings: (i) the \textbf{reference set}, consisting of 70 comments each with more than 200 annotations that served as the foundation for the original MHS construction, and (ii) an \textbf{expanded set} of 5,990 comments, each annotated by at least four human raters where at least 75\% agreed on a single target identity.

\subsection{Language Models}
\label{sec:models}
We consider two distinct sets of models in our analysis. First, the \textbf{primary model set} consists of nine LLMs spanning families, providers, and parameter scales: Claude Opus 4.6, Gemini 3.1 Pro, GPT-5.4, Grok 4.1, DeepSeek V4 Pro, Kimi K2.5, MiniMax M2.5, GPT-OSS 120B, and Llama 3.3 70B. We queried the primary set with both the reference set and the expanded set. We additionally considered an \textbf{expanded model set} of 34 models for which we only queried the reference set. See Appendix~\ref{sec:all_llms} for tables of both primary and expanded model sets with additional metadata. We prompted each LLM with a system prompt that described the task, provided the survey items verbatim from the MHS instrument and specified a JSON response format (see Appendix~\ref{sec:prompts} for prompts). We queried each LLM (default temperature of 1) via its provider's batch API (for closed models) or OpenRouter (for the open-weight models). We used ``medium'' reasoning for all models, where appropriate.

\subsection{Fitting Measurement Models}
\label{sec:measurement_fit}
We fit three classes of measurement models: (i) human-only, (ii) human-anchored, and (iii) LLM-only. The human-only measurement models used the full MHS corpus with a recoding scheme that improved model calibration: \textit{Insult} and \textit{Attack/Defend} were collapsed to four responses, \textit{Humiliate} was collapsed to a ternary label, and \textit{Status}, \textit{Dehumanize}, \textit{Violence}, and \textit{Genocide} were binarized. The human-anchored models had hate speech scores $\theta_n$ and item difficulties $\delta_i$ anchored to those of the human-only measurement scale, with the remaining parameters fit to LLM annotations. The LLM-only measurement model was fit to LLM annotations with no anchoring. We fit all measurement models using the software \textit{Facets} Version 4.5.0 \cite{Linacre2026Facets}.

\section{Results}
\label{sec:results}
We organize our results into two parts. First, we anchor LLM annotations to the measurement scale established by human raters and quantify systematic differences in LLM severity, diagnose the items driving them, and interrogate model bias to question order and target identity (Sections 4.1-4.3). Second, we fit measurement models using LLM annotations alone, diagnose rating scale functioning across items, and compare the resulting LLM scale to its human counterpart (Section 4.4).

\subsection{Measurement Models Quantify Differences between LLM Raters}
We first collected \textbf{reference set} annotations from a \textbf{primary model set} of nine LLMs. As discussed in Section~\ref{sec:measurement_fit}, we fit a measurement model to the LLM annotations that was anchored to the measurement scale established by the human baseline:
\begin{equation}
    \log r_{nijk} = \theta_n^{H} - \delta_i^{H} - \alpha_j^{\text{LLM}} - \tau_{k}\label{eqn:model1}
\end{equation}
where $\theta_n^{H}$ are the hate speech scores and $\delta_i^{H}$ are the item difficulties, both anchored to the values produced by the measurement model fit to human annotations alone. Thus, this model's only degrees of freedom lie in the severity parameters for the LLM raters $\alpha_j^{\text{LLM}}$ and the Rasch-Andrich thresholds $\tau_k$. Fitting Equation~\ref{eqn:model1} to the LLM annotations allows us to precisely answer the question: \textit{given the measurement scale established by human raters, how severe are LLMs when annotating for hate speech?}

We report the fitted severity parameters and compare to the kernel density estimate of human severity values in Figure~\ref{fig:model_severities}a. We find that all LLMs have lower severity values than the median human rater (Fig.~\ref{fig:model_severities}a: negative values). In Equation~\ref{eqn:model1}, lower $\alpha_j$ means the rater flags content more readily than another rater with higher $\alpha_j$. All nine LLMs fall below the median human, indicating that they are more likely to label comments as exhibiting hatefulness aspects.  Since $\alpha_j$ is continuous, we can relate it to external model characteristics. We refit Equation~\ref{eqn:model1} to the \textbf{expanded model set} (Section~\ref{sec:models}) and correlated the $\alpha_j$ with the Artificial Analysis (AA) Intelligence Index (obtained May 2026), a general capability measure \cite{ArtificialAnalysis2026Index}. First, different models from the same provider could have very different severity estimates (Fig.~\ref{fig:model_severities}b: points colored by provider). At the same time, we find that the severity and AA Intelligence Index are significantly positively correlated (Pearson $r = 0.52$, $p < 10^{-3}$), suggesting that more capable models tend to have a higher severity. This pattern could reflect post-training changes over time, the increasing prevalence of reasoning models, or improved ability to make nuanced judgments with higher capability.

%Specific severities are strikingly low: Llama 3.3 70B's severity placed it below the 1st percentile of human raters, ranging up to GPT-5.4's severity in the 38th percentile.

\subsection{Item-Level Interaction Terms Diagnose Miscalibration}
We can extend Equation~\ref{eqn:model1} to incorporate interaction terms that capture the robustness of LLM raters. For example, we demonstrated significant differences between LLM and human raters. We could ask: \textit{what items are driving these differences?} We thus consider the measurement model:
\begin{equation}
\log r_{nijk} = \theta_n^{H} - \delta_i^{H} - \alpha_j^{\text{LLM}} -\tau_k-\mu_{ij} \label{eqn:model2}
\end{equation}
where $\mu_{ij}$ is an interaction term that adjusts the log-odds for rater $j$ and item $i$. We can interpret $\mu_{ij}$ as an item-dependent correction to the severity. Thus, $\mu_{ij}$ serves as a calibration diagnostic: if $\mu_{ij}$ is not significantly different from zero, then item $i$ is well-calibrated for LLM $j$; if it is greater than zero, the LLM is more severe than expected when annotating that item, and vice versa for $\mu_{ij}<0$. 

\begin{figure}[t]
    \centering
    \includegraphics[width=\linewidth]{figures/figure2_severity_decomposition_heatmap.pdf}
    \caption {\textbf{Assessing robustness of item calibration to LLM judges.} Heatmap shows individual $\mu_{ij}$ interaction terms obtained from Equation~\ref{eqn:model2} for items $i$ ($y$-axis) and LLM judges $j$ ($x$-axis). LLMs are sorted from left to right in increasing order of severity (per Fig.~\ref{fig:model_severities}). Color denotes the magnitude and sign of $\mu_{ij}$: red indicates the LLM is less severe than expected (more readily flags), black indicates it is more severe than expected. Significance markers denote Wald $t$-test against $\mu_{ij} = 0$ ($*: p<0.1$; $**: p<0.05$; $***: p<10^{-3}$). }
    \label{fig:severity_decomposition}
\end{figure}

\begin{figure*}[t!]
    \centering
    \includegraphics[width=0.8\linewidth]{figures/figure3_item_order_severity_shift.pdf}
    \caption {\textbf{Quantifying question order bias with DRF.} Models ($y$-axis) are grouped by closed- vs. open-weight (top and bottom, respectively) and sorted alphabetically within each group. \textbf{a.} The original severity $\alpha_j$ calculated from Equation~\ref{eqn:model1} (solid points) and the severity $\alpha_j^R$ calculated on the annotations from the reversed items (open points). \textbf{b.} The change in odds ratio ($x$-axis) induced by the adjustment in severity $\exp(\alpha_j^R - \alpha_j)$. Higher odds ratio indicates that the reversal raises the rating threshold (right side); i.e., the LLM becomes a stricter rater across the items of the construct. Significance markers (right side) denote two-sided Wald $z$-test ($**$: $p<0.05$; $***$: $p<10^{-3}$). Error bars denote one standard error. }
    \label{fig:item_order}
\end{figure*}

% The item-dependent severity adjustments were collectively statistically significant (joint Wald chi-square test: $\chi^2 = 499.6$; $p\ll 10^{-3}$) and explained a minor amount (1 percentage point) of additional variance.

We fit Equation~\ref{eqn:model2} to the primary model annotations and report the individual $\mu_{ij}$ estimates in Figure~\ref{fig:severity_decomposition}. Of the 90 interaction terms ($10$ items $\times$ 9 models), 30 are significant. We identify a few striking trends: first, the interaction term for the \textit{Hate Speech} item is universally negative, suggesting that LLMs are less severe than expected on the specific task of binary hate speech annotation. On the other hand, the \textit{Dehumanize}, \textit{Violence}, and \textit{Genocide} interaction terms were generally positive, suggesting that models were generally more severe -- less likely to give elevated annotations -- than expected. Thus, we can use RMT to diagnose what items drive the differences in severity between LLMs and the human baseline.

%We find them concentrated on the more severe items of the construct.

\subsection{Differential Rater Functioning to Interrogate Model Bias}
RMT provides a suite of techniques called \textit{differential rater functioning} (DRF) which allow a practitioner to assess whether a rater is biased toward specific conditions or sub-groups \cite{linacre1989manyfacet, Maeda2025DIFWords, Minatel2026FairDIF}. A traditional example from the education setting is using DRF to interrogate whether judges are biased toward sub-groups of examinees. Here, we use different DRF techniques to assess two effects: question-order bias and target identity sensitivity.

First, we used the DRF technique of \textit{separated measurements} to assess whether the LLM judges suffered from question order bias. We collected a new set of annotations for the primary model set in which we reversed the question order in the prompt (so that \textit{Genocide} was first and \textit{Sentiment} was last). We fit Equation~\ref{eqn:model1} to both the normal and reversed order, obtaining two severity values per condition, namely $\alpha_{j}$ and $\alpha_{j}^{R}$, respectively. We compare these values in Figure~\ref{fig:item_order}a. We find that, for most models, reversing the question order increases the severity, which suggests that the LLM is less likely to rate comments as exhibiting aspects of hate speech. These effects are only significant for Grok 4.1, MiniMax M2.5, GPT-OSS 120B, and Llama 3.3 70B, the four lowest models by AA Intelligence Index. Thus, significant order effects concentrate among the lowest-capability models. 

This order bias can meaningfully impact annotations. We calculated the change in odds ratio induced by the severity shift, $\exp(\alpha_j^R - \alpha_j)$, which quantifies how much less likely an LLM will raise its rating (e.g., ``agree'' to ``strongly agree'') under the reversed order. We find that odds ratios vary from 1.31 (MiniMax M2.5) to 1.84 (GPT-OSS 120B); the latter suggests GPT-OSS 120B is roughly half as likely to raise its annotation toward the more hateful end of the spectrum under the reversed order as under the original order.

We next considered whether LLMs exhibited different labeling patterns as a function of the \textit{target identity}: the identity group being targeted by a comment. Past work has demonstrated \textit{annotator identity sensitivity}, where raters exhibit different labeling patterns if a comment targets a group they identify with \cite{sachdeva-etal-2022-measuring}. We extend this analysis to LLMs to quantify \textit{target identity sensitivity}. Since the reference set only considers a subset of target identities, we used the primary model set to annotate the \textbf{expanded set} of comments, a sub-corpus of the MHS dataset consisting of 5,990 comments annotated by at least four human raters with high agreement on the target identity (Section~\ref{sec:corpus}). The target identities span four groups (gender, race/ethnicity, religion, and sexuality) and 14 sub-groups.

We fit the measurement model
\begin{equation}
\log r_{nijkm} = \theta_n^{H} - \delta_i^{H} - \alpha_j^{\text{LLM}} -\tau_k-\beta_{jm} \label{eqn:model3}
\end{equation}
which includes a new facet $m$ -- the target identity -- and an interaction term $\beta_{jm}$, describing the adjustment to the severity $\alpha_j^{\text{LLM}}$ needed when a comment targeting identity group $m$ is annotated by judge $j$. We present the odds ratios $\exp(\beta_{jm})$ in Figure~\ref{fig:target_identity_drf}. Odds ratios greater than 1 (gray to black) suggest the LLM judge has a higher ``hate threshold'' for comments targeting that identity group: they are less likely to ascribe hatefulness aspects to comments than what would be expected. Values less than 1 (red) imply a lower hate threshold: when annotating comments targeting this identity group, the LLM judge is more likely to label them as exhibiting aspects of hatefulness.

\begin{figure*}[t]
    \centering
    \includegraphics[width=0.8\linewidth]{figures/figure4_target_drf_heatmap.pdf}
    \caption {\textbf{Interrogating target identity sensitivity with DRF.} Heatmap cells denote the odds ratio corresponding to each interaction term $\exp(\beta_{jm})$ for each LLM judge $j$ ($y$-axis) and target identity group ($x$-axis). Target identities are grouped by category (gender, race/ethnicity, religion, sexuality). Significance markers refer to individual Wald t-tests with null $\exp(\beta_{jm}) = 1$ ($*$: $p<0.1$; $**$: $p<0.05$; $***$: $p<10^{-3}$). }
    \label{fig:target_identity_drf}
\end{figure*}

We find that, on the whole, LLM judges have significant interaction terms (82 of 126 total interaction terms are significant; joint Wald chi-square test: $\chi^2 = 2723; p \ll 10^{-3}$). We identify several notable trends. Comments targeting Muslims universally have odds ratios below 1, implying that LLM judges tend to have lower hate thresholds for these comments. Odds ratios for Jewish-targeting comments are below 1 for most models. Meanwhile, Gay-targeting comments generally elicit higher hate thresholds in LLMs (odds ratios above 1, most significantly so). For race/ethnicity, we observe a split for Black-targeting comments: higher-AA Index models generally have higher hate thresholds, while lower-AA Index models all have lower hate thresholds. Finally, Llama 3.3 70B and Grok 4.1 show the most significant interaction terms overall, indicating they are the most sensitive to target identity. Together, these results demonstrate that RMT can interrogate whether LLM judges exhibit biases to prompt-level conditions (e.g., question order) and content-level attributes (e.g., target identity), particularly on subjective tasks.

\subsection{Constructing an LLM-Only Scale}
Thus far, we have calibrated LLM responses to the measurement scale established by human raters. However, in many deployment settings, LLM judges will set the measurement scale themselves. This raises important questions about whether construct validity and the survey instrument -- often designed with humans in mind -- translate well to the LLM context. Any divergence between LLM- and human-derived scales would mean that the resulting scores carry different meanings than the instrument was designed to measure.

The MHS corpus is well-suited for comparing calibrated LLM and human measurement scales on the same construct. We fit the measurement model
\begin{equation}
\log r_{nijk} = \theta_n - \delta_i - \alpha_j -\tau_k \label{eqn:model4}
\end{equation}
where all of $(\theta_n, \delta_i, \alpha_j, \tau_k)$ are free parameters. We used annotations from the primary model set on the expanded comment set from the MHS corpus (53,910 samples total). We then examined the category functioning ($\tau_k$ values) and compared the LLM and human measurement scales.

% For example, in the MHS corpus, the \textit{Genocide} survey item asked raters whether they agreed that the comment called for genocide of a target identity group using a 5-point Likert scale (``strongly disagree'' to ``strongly agree''). The response thresholds -- $\tau_k$ in Equation~\ref{eqn:model4} -- were found to be miscalibrated, suggesting that human raters were not using the five responses consistently. This motivated collapsing the rating scale from a 5-point Likert to a binary scale: effectively, does the comment call for genocide or not?

\textbf{Diagnosing Disordered Response Thresholds with Rating Scale Functioning.} An often overlooked consideration in evaluations is whether the rating scales associated with each item -- often Likert-like -- are well-constructed. For example, the \textit{Genocide} item in MHS originally had five responses, but was binarized by the corpus authors after the response thresholds were found to be miscalibrated for human raters. Since LLMs have been shown to exhibit systematic biases in their use of Likert rating scales~\cite{tjuatja2024llms, salecha2024llms}, it is worth leveraging RMT to examine whether rating scales translate to the LLM context.

We thus assessed the \textit{rating scale functioning} of each item via the $\tau_k$ thresholds, often referred to as \textit{Rasch-Andrich thresholds} \cite{Andrich1978Rating, Linacre2002Optimizing}. Specifically, each item has five responses (other than the \textit{Hate Speech} item, which has three), corresponding to four (or two) thresholds. A well-calibrated rating scale would always satisfy $\tau_k > \tau_{k-1}$: ``strongly agree'' should have a higher bar than ``agree'' in the log-odds. We plot the Rasch-Andrich thresholds for each item in Figure~\ref{fig:rasch_andrich_thres}. We mark ordered thresholds (where $\tau_k > \tau_{k-1}$) with green arrows pointing right, and disordered thresholds (where $\tau_k < \tau_{k-1}$) with red arrows pointing left. A well-calibrated scale, therefore, would look like three ordered thresholds to the right. We find that some items -- like \textit{Sentiment}, \textit{Respect}, \textit{Humiliate}, and \textit{Attack} -- are well-calibrated. Other rating scales are slightly miscalibrated -- \textit{Insult} and \textit{Status}, for example, only have one disordered threshold, and thus would likely become calibrated after collapsing only one or two adjacent responses. The other three items -- \textit{Dehumanize}, \textit{Violence}, and \textit{Genocide} -- have $\tau_k$ values that generally do not exhibit any global ordering, suggesting they need stronger interventions. Thus, rating scale functioning can facilitate response recoding schemes, with the most severe items in this context requiring the most aggressive collapse.

%Devising a recoding scheme would require leveraging both the observed thresholds and domain knowledge of the survey instrument and construct. In this case, a reasonable scheme would maintain \textit{Sentiment}, \textit{Respect}, \textit{Humiliate}, and \textit{Attack}, collapse \textit{Insult}, \textit{Status} into a ternary scale, and binarize \textit{Dehumanize}, \textit{Violence}, \textit{Genocide}, and \textit{Hate Speech}. Thus, rating scale functioning analysis lets us prescribe a recoding scheme grounded in observed LLM behavior, with the more severe items in the construct requiring the most aggressive collapse.

\begin{figure}[t]
    \centering
    \includegraphics[width=\linewidth]{figures/figure5_category_diagnostics.pdf}
    \caption {\textbf{Category functioning to assess response scale use.} The Rasch-Andrich thresholds $\tau_k$ ($x$-axis) obtained from the LLM measurement scale for each item ($y$-axis). Thresholds are denoted by vertical lines. Ordered steps, where $\tau_k > \tau_{k-1}$, are denoted by green triangles pointing to the right; disordered steps ($\tau_k < \tau_{k-1}$) are denoted by red triangles pointing to the left.}
    \label{fig:rasch_andrich_thres}
\end{figure}


\textbf{Comparing the LLM and Human Measurement Scales.} Finally, we sought to compare the LLM measurement scale to the human measurement scale. To simplify the comparison, we fit Equation~\ref{eqn:model4} to the LLM annotations using the human recoding scheme to ensure both scales operate over identical response categories (see Section~\ref{sec:measurement_fit}) and extracted the comment scores $\theta_n$, item difficulties $\delta_i$, and the nine model severities $\alpha_j$. In general, comparing two measurement scales is not straightforward. We cannot, for example, simply compare the item difficulties between the two scales directly. We can, however, compare \textit{differences} on the logit scale: the difference between, e.g., the \textit{Sentiment} and \textit{Respect} items is comparable between the human and LLM scales. We thus applied a \textit{single-anchor additive linking} constant \cite{Slinde_1979}: we shifted each $\delta_i^H$ by a constant $\Delta = \delta_1^H - \delta_1^{\text{LLM}} $ where $\delta_1$ is the item difficulty for \textit{Sentiment}. We compare the LLM-scale difficulties, (adjusted) human-scale difficulties, the comment scores, and the LLM severities in Figure~\ref{fig:wright_map}. 

We find that the lower end of the human construct -- \textit{Respect}, \textit{Attack}, and \textit{Insult} -- show similar relative spacing across the human and LLM scales. This pattern diverges on the upper end of the construct: \textit{Dehumanize}, \textit{Violence}, and \textit{Genocide} exhibit much greater relative separation, suggesting that LLMs have a much higher threshold for these items. Two items also appear in different rank positions across the two scales: \textit{Hate Speech} occupies a lower-difficulty position on the LLM scale, while \textit{Dehumanize} occupies a higher one. The remaining items preserve their human-scale order. Despite differences in the placement of specific items, the hate speech scores obtained from each measurement scale exhibit high agreement in their order (Spearman correlation $\rho = 0.911$). Together, these results show that the LLM-only scale preserves the ordering of comments on the construct but diverges from the human scale at the upper end, where LLMs require greater endorsement to register severe items -- a finding consistent with the item-level miscalibration identified earlier.

\begin{figure}
    \centering
    \includegraphics[width=\linewidth]{figures/figure6_wright_map.pdf}
    \caption {\textbf{Comparing Human-only and LLM-only Measurement Scales.} Horizontal axis denotes measurement scale for all figure elements. Distribution of comment scores for LLM-only scale denoted by blue curve. Item markers on top half (black) refer to LLM-only difficulties; item markers on bottom half (brown) refer to human-only difficulties with linking to LLM-only scale via \textit{Sentiment}. Model severities denoted by colored points (legend); jitter added for separability.}
    \label{fig:wright_map}
\end{figure}

\section{Discussion}
In this work, we demonstrate how RMT can be an effective tool for interrogating LLMs in the context of evaluation. We focus on the LLM-as-rater paradigm, but evaluation invokes two other paradigms -- LLM-as-examinee and LLM-as-judge -- that would benefit from RMT. LLM-as-examinee benchmarks would improve with RMT's quantification of item difficulty and model ability, exposing whether the benchmark items function as a calibrated instrument. In LLM-as-judge, particularly for subjective tasks, RMT facets can quantify \textit{both} the performance of the LLM examinee and the severity of an LLM judge. 

We conducted a battery of analyses that formalize phenomena previously documented in the evaluation literature \cite{Jiao2026RaterEffects, Wu2026EssayRaters, Jin2025MFRMWriting, Oh2025ChatGPT4oMFRM, Li2025ScoringBias, Xu2026PositionBias}. Severity terms quantify well-documented differences between LLM and human raters on subjective tasks. Differential rater functioning operationalizes an array of LLM biases, such as question order and prompt sensitivity. Rating scale functioning provides a diagnostic for Likert-scale disorder observed in LLM survey responses. RMT unifies these scattered observations into a single quantitative framework anchored to an explicit construct.

A common throughline in our results is a systematic difference between LLM raters and human raters. LLMs have lower severities, treat specific items like \textit{Hate Speech} and \textit{Dehumanize} differently, use rating scale categories
inconsistently on the most severe items, and place items in different positions when constructing their own measurement scale. These findings raise an important point in LLM-as-judge and LLM-as-rater: construct validity established for human raters may not transfer to LLM raters using the same instrument. 

Above all, we emphasize that LLM evaluation should be treated as a measurement problem. We acknowledge that RMT brings substantial overhead -- e.g., construct theorization, instrument design, and validation -- but the field will benefit by adhering to principled design practices. The need is pressing: LLM-as-examinee, -judge, and -rater paradigms are ever more present, and will shape the direction of science, policy, and the field of AI/NLP itself. Measurement theory offers a framework, refined over a century in psychometrics and the social sciences, built for the rigor this problem deserves.


\section{Related Work}
\label{sec:related}
Recent work has argued LLM evaluation is a measurement problem requiring attention to construct validity, scale construction, and score interpretation \cite{Wallach2025GenAIMeasurement, Salaudeen2025MeasurementToMeaning, Zhou2026GeneralScales}. Related IRT and psychometric approaches have examined issues of item difficulty, model ability, and rater effects in NLP/LLM evaluation \cite{lalor2016building, vania2021comparing, rodriguez2021evaluation, Bachmann2024Psychometrics, Federiakin2025Leaderboards, Chen2025CompactIRT, Yao2025JEIRT, Keller2026StatisticalModels, Choi2026JudgeIRT, Feng2026NoisyJudges, Chen2026NoisyJudgeInference, Jiao2026RaterEffects, Wu2026EssayRaters, Jin2025MFRMWriting}. We build on this work by centering Rasch measurement theory, a related approach. For hate speech and other data perspectivist tasks \cite{cabitza2023toward}, LLMs are increasingly used as scalable raters; we directly contribute to this literature by shedding light on how modern LLMs behave for these tasks \cite{sap2022annotators,giorgi2025human, Okpala2025LLMAnnotationBias}.

\section*{Limitations}

While we broadly aimed to convey that RMT is relevant for a variety of measurement tasks useful for LLM evaluation, our specific analysis is restricted to hate speech via the MHS corpus. We do not empirically validate that RMT can easily be applied in other contexts -- indeed, this will likely require additional construct and instrument development, or at least adaptation of existing benchmarks within a viable RMT framework. In the same vein, we propose that RMT is useful for LLM-as-judge and LLM-as-examinee, but do not demonstrate this in specific datasets, instead opting to spend the space of this paper on depth of analysis for a single use case.

Our analysis is restricted to the MHS corpus, which is largely in English. Thus, for purposes of contextualizing our results in the hate speech literature, our results should not be interpreted in a multilingual context. Additionally, our analyses were limited to one sample per comment. We did, however, consider a large amount of comments: up to 5,990 for specific analyses, so it is unlikely the standard errors increase much more when including variation from generation stochasticity. Similarly, while we considered sensitivity to question order, it is possible that other prompt constructions could impact the parameter estimates we obtain.


\section*{Acknowledgments}
We used AI Assistants to support the analyses conducted in this paper:
\begin{itemize}
    \item We used agentic coding tools to develop most of the code use to query LLMs, analyze data, and create plots shown in this paper, under guidance of the authors.
    \item We used LLMs to assist with literature review.
    \item We used LLMs to provide feedback on wording, tone, and grammar (but all writing was done by the authors).
\end{itemize}

\section*{Ethical Considerations}
We used the \textit{Measuring Hate Speech} corpus, a dataset of social media comments often containing offensive language and slurs. We used the dataset in compliance with the authors' intended use. In particular, personal identifying information were removed from all comments. 

Our analyses raise the possibility of LLMs replacing humans for annotation tasks. For sensitive topics, like hate speech detection, replacement of human content moderators may raise questions of adverse and unintended consequences, particularly on specific identity groups. Our results are not intended to advocate for replacement of humans for such tasks; instead, we intend for our results to demonstrate that such tasks may \textit{not} transfer well even if some benchmarks suggest that LLMs are highly ``accurate.'' 

% Bibliography entries for the entire Anthology, followed by custom entries
%\bibliography{custom,anthology-overleaf-1,anthology-overleaf-2}

% Custom bibliography entries only
\bibliography{custom}

\clearpage
\appendix

\section{System Prompts}
\label{sec:prompts}
\begin{tcblisting}{
  enhanced,
  breakable,
  colback=white,
  colframe=black,
  title={System Prompt},
  fonttitle=\bfseries,
  coltitle=white,
  width=\linewidth,
  boxsep=3pt,
  left=4pt,
  right=4pt,
  top=4pt,
  bottom=4pt,
  before skip=6pt,
  after skip=6pt,
  listing only,
  listing engine=listings,
  listing options={
    basicstyle=\ttfamily\scriptsize,
    breaklines=true,
    breakatwhitespace=false,
    columns=fullflexible,
    keepspaces=true,
    showstringspaces=false,
    aboveskip=0pt,
    belowskip=0pt
  }
}
You are a tool for identifying and characterizing hate speech.

INSTRUCTIONS:
- Analyze the provided social media comment using the target-group question and the 10 survey items below.
- The comment may or may not contain hate speech.
- Answer every field using only the allowed lettered response options.
- For target_groups, return one or more uppercase letters in a JSON array.
- For sentiment, respect, insult, humiliate, status, dehumanize, violence, genocide, attack_defend, and hate_speech, return exactly one uppercase letter.
- Respond only with a single valid JSON object matching the required schema.
- Do not include explanations, summaries, markdown, or extra text.

TARGET-GROUP QUESTION:
1. Is the provided comment directed at or about any individual or groups based on: select one or more

A. Race or ethnicity
B. Religion
C. National origin or citizenship status
D. Gender
E. Sexual orientation
F. Age
G. Disability status
H. Political ideology
I. None of the above

SURVEY ITEMS:
2. How would you describe the sentiment of this comment?
A. Strongly negative
B. Somewhat negative
C. Neutral
D. Somewhat positive
E. Strongly positive

3. Is this comment respectful towards the group(s) you previously identified?
A. Strongly disrespectful
B. Disrespectful
C. Neutral
D. Respectful
E. Strongly respectful

4. Do you agree or disagree with the following statement? This comment is insulting towards the group(s) you previously identified.
A. Strongly disagree
B. Disagree
C. Neither disagree nor agree
D. Agree
E. Strongly agree

5. Do you agree or disagree with the following statement? This comment is humiliating towards the group(s) you previously identified.
A. Strongly disagree
B. Disagree
C. Neither disagree nor agree
D. Agree
E. Strongly agree

6. This comment states that the group(s) you previously identified is:
A. Strongly inferior
B. Inferior
C. Neither superior nor inferior
D. Superior
E. Strongly superior

7. Do you agree or disagree with the following statement? This comment dehumanizes the group(s) you previously identified (e.g., by comparing them to an animal).
A. Strongly disagree
B. Disagree
C. Neither disagree nor agree
D. Agree
E. Strongly agree

8. Do you agree or disagree with the following statement? This comment calls for using violence against the group(s) you previously identified.
A. Strongly disagree
B. Disagree
C. Neither disagree nor agree
D. Agree
E. Strongly agree

9. Do you agree or disagree with the following statement? This comment calls for the deliberate killing of a large group of people from the group(s) you previously identified.
A. Strongly disagree
B. Disagree
C. Neither disagree nor agree
D. Agree
E. Strongly agree

10. Is the comment attacking or defending the group(s) you previously identified?
A. Strongly defending
B. Defending
C. Neither defending nor attacking
D. Attacking
E. Strongly attacking

11. Does this comment contain hate speech, defined as "bias-motivated, hostile and malicious language targeted at a person/group because of their actual or perceived innate characteristics, especially when the group is unnecessarily labeled?"
A. Yes
B. No
C. Unclear

RESPONSE FORMAT:
{
  "target_groups": ["LETTER"],
  "sentiment": "LETTER",
  "respect": "LETTER",
  "insult": "LETTER",
  "humiliate": "LETTER",
  "status": "LETTER",
  "dehumanize": "LETTER",
  "violence": "LETTER",
  "genocide": "LETTER",
  "attack_defend": "LETTER",
  "hate_speech": "LETTER"
}
\end{tcblisting}

\captionof{figure}{System prompt used for model annotation.}
\label{fig:system-prompt}

\begin{tcblisting}{
  enhanced,
  breakable,
  colback=white,
  colframe=black,
  title={System Prompt},
  fonttitle=\bfseries,
  coltitle=white,
  width=\linewidth,
  boxsep=3pt,
  left=4pt,
  right=4pt,
  top=4pt,
  bottom=4pt,
  before skip=6pt,
  after skip=6pt,
  listing only,
  listing engine=listings,
  listing options={
    basicstyle=\ttfamily\scriptsize,
    breaklines=true,
    breakatwhitespace=false,
    columns=fullflexible,
    keepspaces=true,
    showstringspaces=false,
    aboveskip=0pt,
    belowskip=0pt
  }
}
You are a tool for identifying and characterizing hate speech.

INSTRUCTIONS:
- Analyze the provided social media comment using the target-group question and the 10 survey items below.
- The comment may or may not contain hate speech.
- Answer every field using only the allowed lettered response options.
- For target_groups, return one or more uppercase letters in a JSON array.
- For sentiment, respect, insult, humiliate, status, dehumanize, violence, genocide, attack_defend, and hate_speech, return exactly one uppercase letter.
- Respond only with a single valid JSON object matching the required schema.
- Do not include explanations, summaries, markdown, or extra text.

TARGET-GROUP QUESTION:
1. Is the provided comment directed at or about any individual or groups based on: (select one or more)
A. Race or ethnicity
B. Religion
C. National origin or citizenship status
D. Gender
E. Sexual orientation
F. Age
G. Disability status
H. Political ideology
I. None of the above

SURVEY ITEMS:
2. Do you agree or disagree with the following statement: This comment calls for the deliberate killing of a large group of people from the group(s) you previously identified.
A. Strongly disagree
B. Disagree
C. Neither disagree nor agree
D. Agree
E. Strongly agree

3. Do you agree or disagree with the following statement: This comment calls for using violence against the group(s) you previously identified.
A. Strongly disagree
B. Disagree
C. Neither disagree nor agree
D. Agree
E. Strongly agree

4. Do you agree or disagree with the following statement: This comment dehumanizes the group(s) you previously identified.
A. Strongly disagree
B. Disagree
C. Neither disagree nor agree
D. Agree
E. Strongly agree

5. This comment states that the group(s) you previously identified is:
A. Strongly inferior
B. Inferior
C. Neither superior nor inferior
D. Superior
E. Strongly superior

6. Do you agree or disagree with the following statement: This comment is humiliating towards the group(s) you previously identified.
A. Strongly disagree
B. Disagree
C. Neither disagree nor agree
D. Agree
E. Strongly agree

7. Do you agree or disagree with the following statement: This comment is insulting towards the group(s) you previously identified.
A. Strongly disagree
B. Disagree
C. Neither disagree nor agree
D. Agree
E. Strongly agree

8. Is this comment respectful towards the group(s) you previously identified?
A. Strongly disrespectful
B. Disrespectful
C. Neutral
D. Respectful
E. Strongly respectful

9. How would you describe the sentiment of this comment?
A. Strongly negative
B. Somewhat negative
C. Neutral
D. Somewhat positive
E. Strongly positive

10. Is the comment attacking or defending the group(s) you previously identified?
A. Strongly defending
B. Defending
C. Neither defending nor attacking
D. Attacking
E. Strongly attacking

11. Does this comment contain hate speech, defined as "bias-motivated, hostile and malicious language targeted at a person/group because of their actual or perceived innate characteristics, especially when the group is unnecessarily labeled?"
A. Yes
B. No
C. Unclear

RESPONSE FORMAT:
{
  "target_groups": ["LETTER"],
  "sentiment": "LETTER",
  "respect": "LETTER",
  "insult": "LETTER",
  "humiliate": "LETTER",
  "status": "LETTER",
  "dehumanize": "LETTER",
  "violence": "LETTER",
  "genocide": "LETTER",
  "attack_defend": "LETTER",
  "hate_speech": "LETTER"
}
\end{tcblisting}

\captionof{figure}{System prompt used for model annotation.}
\label{fig:system-prompt-reversed}


\section{Exhaustive List of LLMs}
\label{sec:all_llms}

See Table 1 for the primary model set and Table 2 for the expanded model set.

\begin{table*}[t!]
\centering
\caption{Primary model set. Refusal rate is computed over the 5,990 full-set MHS comments.}
\label{tab:primary-models}
\begin{tabular}{llllr}
\toprule
Provider & Model & Access & AA Index & Refusal Rate \\
\midrule
OpenAI & GPT-5.4 & Closed & 47.9 & 0.0\% (0/5{,}990) \\
Google & Gemini 3.1 Pro Preview & Closed & 57.2 & 0.7\% (44/5{,}990) \\
Anthropic & Claude Opus 4.6 & Closed & 46.5 & 0.0\% (0/5{,}990) \\
xAI & Grok 4.1 Fast Reasoning & Closed & 38.6 & 0.0\% (0/5{,}990) \\
DeepSeek & DeepSeek V4 Pro & Open-Weights & 51.5 & 0.0\% (0/5{,}990) \\
MiniMax & MiniMax M2.5 & Open-Weights & 41.9 & 0.0\% (0/5{,}990) \\
Moonshot AI & Kimi K2.5 & Open-Weights & 46.8 & 1.7\% (104/5{,}990) \\
OpenAI & GPT-OSS 120B & Open-Weights & 33.3 & 0.0\% (0/5{,}990) \\
Meta & Llama 3.3 70B Instruct Turbo & Open-Weights & 14.5 & 0.0\% (0/5{,}990) \\
\bottomrule
\end{tabular}
\end{table*}

\clearpage
\begin{table*}[t!]
\centering
\caption{Expanded Figure 1 scatter-plot model set. Refusal rate is computed over the 70 MHS reference-set comments.}
\label{tab:expanded-scatter-models}
\begin{tabular}{llllr}
\toprule
Provider & Model & Access & AA Index & Refusal Rate \\
\midrule
OpenAI & GPT-4.1 & Closed & 26.3 & 0.0\% (0/70) \\
OpenAI & GPT-4o & Closed & 17.3 & 0.0\% (0/70) \\
OpenAI & GPT-5.2 & Closed & 46.6 & 0.0\% (0/70) \\
OpenAI & GPT-5.4 & Closed & 47.9 & 0.0\% (0/70) \\
OpenAI & GPT-5.4 Mini & Closed & 37.7 & 0.0\% (0/70) \\
OpenAI & GPT-5.4 Nano & Closed & 38.1 & 0.0\% (0/70) \\
Google & Gemini 2.5 Flash & Closed & 20.6 & 1.4\% (1/70) \\
Google & Gemini 2.5 Pro & Closed & 34.6 & 2.9\% (2/70) \\
Google & Gemini 3 Flash Preview & Closed & 46.4 & 0.0\% (0/70) \\
Google & Gemini 3.1 Flash Lite Preview & Closed & 33.5 & 0.0\% (0/70) \\
Google & Gemini 3.1 Pro Preview & Closed & 57.2 & 0.0\% (0/70) \\
Anthropic & Claude Haiku 4.5 & Closed & 31.0 & 0.0\% (0/70) \\
Anthropic & Claude Sonnet 4 & Closed & 33.0 & 11.4\% (8/70) \\
Anthropic & Claude Sonnet 4.5 & Closed & 37.1 & 1.4\% (1/70) \\
Anthropic & Claude Sonnet 4.6 & Closed & 44.4 & 0.0\% (0/70) \\
Anthropic & Claude Opus 4 & Closed & 33.0 & 2.9\% (2/70) \\
Anthropic & Claude Opus 4.1 & Closed & 36.0 & 2.9\% (2/70) \\
Anthropic & Claude Opus 4.5 & Closed & 49.7 & 0.0\% (0/70) \\
Anthropic & Claude Opus 4.6 & Closed & 46.5 & 0.0\% (0/70) \\
xAI & Grok 3 & Closed & 25.2 & 0.0\% (0/70) \\
xAI & Grok 4 Fast Non-Reasoning & Closed & 23.1 & 0.0\% (0/70) \\
xAI & Grok 4 Fast Reasoning & Closed & 35.1 & 0.0\% (0/70) \\
xAI & Grok 4.1 Fast Non-Reasoning & Closed & 23.6 & 0.0\% (0/70) \\
xAI & Grok 4.1 Fast Reasoning & Closed & 38.6 & 0.0\% (0/70) \\
DeepSeek & DeepSeek V3.2 & Open-Weights & 32.1 & 0.0\% (0/70) \\
DeepSeek & DeepSeek V4 Pro & Open-Weights & 51.5 & 0.0\% (0/70) \\
MiniMax & MiniMax M2.5 & Open-Weights & 41.9 & 0.0\% (0/70) \\
Moonshot AI & Kimi K2.5 (direct) & Open-Weights & 46.8 & 7.1\% (5/70) \\
Moonshot AI & Kimi K2.5 (OpenRouter) & Open-Weights & 46.8 & 0.0\% (0/70) \\
OpenAI & GPT-OSS 120B & Open-Weights & 33.3 & 0.0\% (0/70) \\
Meta & Llama 3.3 70B Instruct Turbo & Open-Weights & 14.5 & 0.0\% (0/70) \\
Qwen & Qwen3.5 122B A10B & Open-Weights & 41.6 & 0.0\% (0/70) \\
Xiaomi & MiMo V2 Pro & Open-Weights & 49.2 & 0.0\% (0/70) \\
Z.ai & GLM-5 Turbo & Closed & 46.8 & 0.0\% (0/70) \\
\bottomrule
\end{tabular}
\end{table*}


\end{document}
